\usepackage[utf8]{inputenc}
\usepackage[T1]{fontenc}
\usepackage{lmodern}
\usepackage{amsthm} %ou \usepackage{ntheorem}
\usepackage[top=1.5cm, bottom=2.2cm, left=2cm, right=2cm]{geometry}
\usepackage{fancyhdr}
\usepackage{graphicx}
\usepackage{draftwatermark}
\usepackage{multicol}
\usepackage{enumerate}
\usepackage{titlesec}
\usepackage{alltt}
\usepackage{kpfonts}
\usepackage[svgnames,table]{xcolor}
\usepackage[francais]{babel}
\pagestyle{fancy}


%
% Création du filigrane
%
\SetWatermarkText{\includegraphics{../logo_blanc.png}}
\SetWatermarkScale{0.5}
\SetWatermarkAngle{0}


%
% Suivent les macros qui donnent la mise en forme des théorèmes, définitions, etc...
%
\newtheoremstyle{persoth}% name
{2pt}%Space above
{2pt}%Space below
{}%{\itshape}%Body font
{}%Indent amount
{\bf}%Theorem head font
{.}%Punctuation after theorem head
{.5em}%Space after theorem head 2
{}%

\newtheoremstyle{persodef}% name
{2pt}%Space above
{2pt}%Space below
{}%Body font
{}%Indent amount
{\bf}%Theorem head font
{.}%Punctuation after theorem head
{.5em}%Space after theorem head 2
{}%

\theoremstyle{persoth}% default
\newtheorem*{thm}{\noindent\textcolor{Crimson}{Th}}
\newtheorem*{lem}{\noindent\textcolor{MediumVioletRed}{Lem}}
\newtheorem*{pop}{\noindent\textcolor{FireBrick}{Prop}}
\newtheorem*{cor}{\noindent\textcolor{Brown}{Cor}}

\theoremstyle{persodef}
\newtheorem*{defn}{\noindent\textcolor{magenta}{Def}}
\newtheorem*{conj}{Conjecture}

\theoremstyle{remark}
\newtheorem*{rem}{\noindent\textcolor{Teal}{Rem}}
\newtheorem*{ex}{\noindent\textcolor{DarkOrange}{Ex}}
\newtheorem*{note}{\noindent\textcolor{RoyalBlue}{Not}}
\newtheorem*{danger}{\textcolor{green}{Attention}}
\newtheorem*{voc}{\textcolor{DarkGreen}{Voc}}
\newtheorem*{hyp}{\textcolor{OrangeRed}{Hyp}}

%\newcommand\demo{\begin{proof}[\textit{Démonstration}]}


%
% Redéfinition de caractéristiques sur le format des titres de sections etc...
%
\titleformat{\section}
	[hang]% style : hang, display, runin, leftmargin, ...
	{\Large\bfseries}% fonte numéro + titre
	{\thesection}% numéro
	{1em}% espace entre le numéro et le titre
	{}% fonte titre

\titlespacing{\section}{0pt}
	{8pt}{3pt}

\titleformat{\subsection}
	[hang]% style : hang, display, runin, leftmargin, ...
	{\normalsize\bfseries\sffamily}% fonte numéro + titre
	{\thesubsection}% numéro
	{4pt}% espace entre le numéro et le titre
	{}% fonte titre

\titlespacing{\subsection}{0pt}
	{6pt}{1pt}


\date{}

%
% Ce qui suit permet de corriger le bug de titlesec sur la numérotation des sections
%
\usepackage{etoolbox}
\makeatletter
\patchcmd{\ttlh@hang}{\parindent\z@}{\parindent\z@\leavevmode}{}{}
\patchcmd{\ttlh@hang}{\noindent}{}{}{}
\makeatother
